\documentclass{tpks}

\doctitle{План управления конфигурацией ПО}
\docversion{1}
\docdate{17.10.2026}

\begin{document}
\makeheader

\section{Введение}

\subsection{Цель}

План задаёт порядок управления конфигурацией программного обеспечения проекта. УК ПО
должно сохранять целостность данных жизненного цикла: участники должны работать с
известными версиями, изменения должны иметь основание и результат, а утверждённую
конфигурацию должно быть возможно воспроизвести.

\subsection{Область применения}

План действует для документов, требований, исходного кода, сценариев проверки,
результатов сборки и записей об их изменениях. Он применяется от начала подготовки
изменения до выпуска и загрузки утверждённой конфигурации. Объектный код,
полученный компилятором, не редактируется и не является самостоятельной ЕК, но его
версия и контрольная сумма фиксируются в протоколе интеграции.

\subsection{Ссылки}

\begin{itemize}
  \item DO-178C и КТ-178С, раздел 7 «Управление конфигурацией ПО»;
  \item системные требования проекта, версия 3;
  \item стандарт на разработку требований к ПО проекта (после утверждения).
\end{itemize}

\subsection{Термины и сокращения}

\printterms{Базовая версия, ЕК, ГУИ, ЗИ, КК1, КК2, СКУ, СП, УК ПО}

В этом плане используются также термины GitHub, issue, pull request, ветка, тег,
релиз и артефакт сборки. Их значения определены контекстом GitHub и правилами
данного плана. Issue используется для записи СП или ЗИ, но тип записи всегда
указывается меткой.

\section{Среда управления конфигурацией ПО}

\subsection{Инструменты УК ПО}

Единственным разрешённым источником данных проекта служит репозиторий команды на
GitHub. В нём хранятся документы в каталоге \texttt{docs}, готовые PDF в
\texttt{pdf}, исходный код, сценарии проверки, файлы сборки и записи процесса.
Папка \texttt{materials} содержит исходные учебные материалы и не изменяется в
рамках разработки.

\begin{tabl}{|l|L|}{Инструменты и назначение}[tab:scm-tools]{\textbf{Инструмент} & \textbf{Назначение}}
Git & Хранение версий файлов, веток, коммитов и тегов. \\ \hline
GitHub Issues & Регистрация СП и ЗИ, их статусы, обсуждение и трассировка. \\ \hline
GitHub Pull requests & Рассмотрение изменений, обсуждение строк и разрешение слияния. \\ \hline
GitHub Actions & Автоматическая сборка и запуск проверок после создания pull request и после слияния. \\ \hline
Хранилище артефактов & Хранение утверждённых PDF, протоколов сборки и выпусков. \\ \hline
\end{tabl}

Доступ к репозиторию предоставляется только участникам команды. Право чтения
нужно для просмотра конфигурации и проверки результатов, право записи --- для
создания веток, коммитов и pull request. Слияние в \texttt{main} разрешено только
проверяющему изменения или ответственному за выпуск после выполнения критериев
перехода из раздела 5.

\clearpage
\subsection{Организация и ответственность}

\begin{tabl}{|l|L|}{Роли процесса УК ПО}[tab:scm-roles]{\textbf{Роль} & \textbf{Ответственность}}
Совет по управлению изменениями & Коллегиально решает по СП, нужны ли изменения и какие ЕК затрагиваются. \\ \hline
ГУИ & Постоянная часть совета; организует рассмотрение и контролирует состояние СП и ЗИ. \\ \hline
Руководитель проекта & Назначает автора ЗИ и проверяющего, разрешает спорные организационные вопросы. \\ \hline
Автор изменений & Вносит изменение только в пределах назначенного ЗИ, обновляет трассировку и описание результата. \\ \hline
Проверяющий изменения & Проверяет соответствие изменения ЗИ и входным данным, анализирует результаты CI и принимает решение по pull request. \\ \hline
Ответственный за выпуск & Формирует выпуск из утверждённых ЕК, публикует его и контролирует загрузку. \\ \hline
\end{tabl}

Автор изменения не может быть единственным проверяющим этого изменения. Состав
ГУИ и правило назначения МИ, указанного на схеме жизненного цикла ЗИ,
нужно согласовать с командой.

\subsection{Единицы конфигурации и категории контроля}

ЕК однозначно определяется парой «уникальное имя, версия». Составная ЕК имеет
собственную версию и содержит версии входящих в неё ЕК. Версия ЕК увеличивается
только по утверждённому ЗИ.

\clearpage
\begin{tabl}{|l|l|L|}{Единицы конфигурации}[tab:configuration-items]{\textbf{ЕК} & \textbf{Категория} & \textbf{Идентификация и хранение}}
Системные требования и СПТ & КК1 & Файл с именем документа и версией; исходный текст и PDF в Git. \\ \hline
Планы и стандарты & КК1 & Файл документа с номером версии; исходный текст и PDF в Git. \\ \hline
Исходный код и сценарии проверки & КК1 & Путь файла, коммит и версия репозитория. \\ \hline
СП и ЗИ & КК1 & Issue с уникальным номером, типом и списком затронутых ЕК. \\ \hline
Pull request & КК2 & Номер репозитория и PR, связанные с ЗИ и коммитами. \\ \hline
Протоколы ФИ, сборки и загрузки & КК2 & Файл протокола с датой, версиями входов и результатом. \\ \hline
Конфигурация CI и зависимости & КК1 & Файл конфигурации, коммит и зафиксированные версии зависимостей. \\ \hline
Выпуск ПО и его манифест & КК2 & Тег релиза, список ЕК и контрольные суммы артефактов. \\ \hline
\end{tabl}

КК1 означает, что после создания базовой версии изменение допускается только по
ЗИ. Для КК2 достаточно фиксировать версию, автора, дату и результат операции;
порядок повторной проверки задаётся видом протокола. Объектный код не
редактируется и поэтому не получает отдельную категорию контроля.

\subsection{Средства трассировки}

Каждая запись СП перечисляет ЕК, в которых обнаружена проблема. Утверждённое СП
связано с одним или несколькими ЗИ, а каждый ЗИ --- с одной изменяемой ЕК,
веткой, pull request и итоговым коммитом. Манифест выпуска связывает версии всех
включённых ЕК с артефактами и протоколом загрузки. Связь между версиями одной ЕК
проходит через ЗИ; промежуточные коммиты не считаются базовыми версиями.

\subsection{Связь с другими процессами ЖЦ ПО}

Разработка требований, проектирование и кодирование получают ЗИ как основание
изменения и передают результаты в УК ПО для проверки и фиксации. Интеграция
создаёт исполняемый объектный код, протокол сборки и протокол загрузки. Верификация
и гарантия качества регистрируют найденные несоответствия через СП. Планирование
определяет инструменты и правила, которые далее контролируются этим планом.

\section{Мероприятия процесса УК ПО}

\subsection{Идентификация конфигурации}

Перед началом работы ГУИ определяет состав данных жизненного цикла и их связи.
Для каждой ЕК задаются уникальное имя, расположение, категория контроля и правило
версии. Идентификатор документа состоит из имени файла и номера версии; исходный
код идентифицируется путём и коммитом; issue --- номером и меткой; pull request
--- номером PR; выпуск --- тегом.

Ветка \texttt{main} содержит только проверенные изменения. Рабочая ветка
создаётся для одного ЗИ и именуется \texttt{zi-\textless номер\textgreater
-\textless краткое-описание\textgreater}. Коммиты должны содержать номер ЗИ,
например \texttt{ZI\#17: update SCM plan}. Удаление коммитов и истории после
слияния запрещено.

\subsection{Базовые версии и трассируемость}

\subsubsection{Создание базовой версии}

Рабочие версии могут существовать в ветке ЗИ, но базовой становится только версия,
которая прошла проверку, утверждена и передана в СКУ. Для документов базовая
версия фиксируется коммитом в \texttt{main}, а для выпуска ПО дополнительно тегом
вида \texttt{v\allowbreak\textless major\textgreater.\allowbreak\textless
minor\textgreater.\allowbreak\textless patch\textgreater}. Формулировка «самый
новый коммит \texttt{main}» означает только последний коммит, вошедший в
\texttt{main} по
утверждённому ЗИ; она не делает любой коммит автоматически базовым.

\subsubsection{Контроль библиотек}

Файлы зависимостей входят в КК1. В них фиксируются версии библиотек и контрольные
суммы или lock-файлы. CI устанавливает зависимости только из разрешённого
источника и сохраняет протокол установки. Обновление зависимости выполняется
отдельным ЗИ с анализом влияния.

\subsubsection{Трассируемость версий}

Перед слиянием проверяющий убеждается, что PR связан с ЗИ, ЗИ связано с
утверждённым СП, а список изменённых файлов совпадает со списком ЕК в ЗИ. После
слияния в ЗИ записываются хэш коммита, версии документов и ссылка на протокол CI.

\subsection{Регистрация и отслеживание проблем}

СП регистрирует одно несоответствие между входами и выходами процессов, данными
жизненного цикла и стандартами либо процессами и планами. Зарегистрировать СП
может любой участник проекта.

В СП указываются уникальный номер, краткое описание проблемы, шаги
воспроизведения, фактический и ожидаемый результат, список затронутых ЕК и их
версий. СП описывает проблему, а не предполагаемое решение. Для записи
используется метка \texttt{type::problem}.

Жизненный цикл СП:
\begin{enumerate}
  \item «Создано» --- запись ещё не проверена;
  \item «Зарегистрировано» --- ГУИ проверила обязательные поля;
  \item «Рассматривается» --- совет оценивает проблему и влияние;
  \item «В работе» --- по СП выполняются утверждённые изменения;
  \item «Завершено» --- изменения внесены, проверены и связаны с результатом;
  \item «Отклонено» --- запись создана ошибочно или изменений не требуется.
\end{enumerate}

\subsection{Управление изменениями: СП и ЗИ}

Совет по управлению изменениями рассматривает зарегистрированное СП. Он может
отклонить СП, если проблема подтверждена ошибочно или изменение не требуется,
либо утвердить его и перечислить затрагиваемые ЕК. Из одного СП порождается
несколько ЗИ, по одному на каждую изменяемую ЕК.

ЗИ описывает планируемое изменение, его основание, исходную версию ЕК, критерий
готовности и способ проверки. Для него используется метка \texttt{type::change}.
Жизненный цикл ЗИ:
\begin{enumerate}
  \item «Создан» --- ЗИ создан из утверждённого СП;
  \item «Передан в работу» --- руководитель проекта назначил автора;
  \item «В работе» --- автор меняет ЕК в рабочей ветке;
  \item «Проверка» --- изменение готово, проверяющий выполняет ревью и проверки;
  \item «Завершён» --- изменение утверждено, слито и зафиксировано;
  \item «Отклонён» --- изменение ошибочно или больше не требуется.
\end{enumerate}

\subsection{Рассмотрение изменений}

Автор создаёт ветку ЗИ от актуального \texttt{main}, вносит только относящиеся к
ЗИ коммиты и открывает pull request. В описание PR он добавляет ссылки на СП и
ЗИ, перечень ЕК, описание проверки и ожидаемый результат. CI запускает сборку,
проверку формата, тесты и проверку ссылок.

Проверяющий сравнивает изменение с ЗИ, проверяет трассируемость, результаты CI и
отсутствие открытых замечаний. Если есть несоответствие, PR возвращается автору,
а ЗИ остаётся в состоянии «В работе». После устранения замечаний проверяющий
повторяет проверку. Только после одобрения PR ответственный за слияние выполняет
merge в \texttt{main}; удалять PR и историю обсуждения запрещено.

\subsection{Интеграция изменений}

После слияния CI выполняет интеграционную сборку из коммита \texttt{main}.
Входами являются исходный код, параметры сборки, описание проекта, компонент
параметрических данных (если он есть), План разработки ПО и ЗИ. Выходами
являются исполняемый объектный код, протокол сборки, командные файлы сборки,
протокол загрузки и сообщения о проблемах.

Артефакты и протоколы получают идентификаторы, связываются с коммитом \texttt{main}
и сохраняются в хранилище. Если сборка или загрузка завершилась ошибкой, CI
создаёт СП с указанием версии входов. Объектный код не изменяется вручную и не
используется как источник для следующей сборки.

\subsection{Хранение, воспроизведение и выпуск}

GitHub защищает \texttt{main} от прямой записи и хранит историю коммитов, issue и
PR. Резервная копия репозитория и хранилища артефактов выполняется регулярно;
ответственный за УК проверяет возможность чтения копии. Воспроизведение
конфигурации выполняется по тегу выпуска, манифесту ЕК, коммиту \texttt{main} и
протоколу сборки.

Ответственный за выпуск проверяет, что все ЕК выпуска имеют утверждённые версии,
PR слиты, СП и ЗИ закрыты, а протокол интеграции не содержит открытых
несоответствий. Затем он создаёт тег, формирует манифест и публикует выпуск в
хранилище. Непроверенные или изменённые вручную файлы в выпуск не включаются.

\subsection{Контроль загрузки ПО}

Перед загрузкой ответственный за выпуск сверяет тег, манифест и контрольные
суммы артефактов. Загрузку выполняют только из опубликованного выпуска. После
загрузки фиксируются целевая среда, версия аппаратуры, время, исполнитель,
контрольная сумма и результат проверки. Несовпадение контрольной суммы или
несовместимость с аппаратурой блокирует загрузку и регистрируется СП.

\subsection{Контроль среды жизненного цикла ПО}

Версии Git, среды сборки, инструментов проверки, зависимостей и раннеров CI
перечисляются в файле конфигурации проекта. Этот файл является ЕК КК1. Изменение
инструмента или его версии выполняется по ЗИ с проверкой воспроизводимости
сборки. Доступы к GitHub и хранилищу выдаются по роли и отзываются при уходе
участника из проекта.

\subsection{Учёт состояния конфигурации}

Состояние конфигурации отражается в GitHub и в манифесте выпуска. Для каждой ЕК
доступны её версия, категория, владелец изменения, связанное СП или ЗИ, состояние
проверки и базовая версия. ГУИ перед выпуском сверяет таблицу состояний со
списком файлов, PR, протоколами и тегом.

\section{Документы управления конфигурацией}

К документам УК ПО относятся настоящий план, шаблоны СП и ЗИ, правила веток и
коммитов, матрица ЕК, манифесты выпусков, протоколы CI, сборки и загрузки,
резервные копии и отчёты о состоянии конфигурации. У каждого документа есть
версия и запись в листе изменений. Утверждённые документы хранятся в
\texttt{docs} и PDF в \texttt{pdf}; рабочие черновики не являются базовыми
версиями.

\section{Критерии перехода}

\clearpage
\begin{tabl}{|l|L|}{Критерии перехода}[tab:transition-criteria]{\textbf{Переход} & \textbf{Условие}}
СП создано $\rightarrow$ зарегистрировано & Заполнены идентификатор, одна проблема, затронутые ЕК и их версии. \\ \hline
СП рассматривается $\rightarrow$ утверждено или отклонено & Совет зафиксировал решение и список ЕК. \\ \hline
ЗИ создано $\rightarrow$ передано в работу & Назначены автор и проверяющий, указаны критерии готовности. \\ \hline
ЗИ в работе $\rightarrow$ проверка & PR открыт, CI завершён успешно, трассировка заполнена. \\ \hline
ЗИ проверка $\rightarrow$ завершён & Проверяющий одобрил PR, изменение слито в \texttt{main}, протокол сохранён. \\ \hline
Конфигурация $\rightarrow$ выпуск & Все ЕК утверждены, интеграция успешна, манифест и контрольные суммы сформированы. \\ \hline
Выпуск $\rightarrow$ загрузка & Ответственный сверил тег, манифест, артефакты и совместимость целевой среды. \\ \hline
\end{tabl}

\begin{changelog}
\change{1}{17.10.2026}{Первая версия плана управления конфигурацией ПО}{Лазарев Я. О.}{Задание к паре 17.10.2026}
\end{changelog}

\end{document}
